% ===========================================================================
% Experiment 02: the DIT placement crossover, side by side.
%
% OVERLEAF: upload this file and the three PDFs beside it into a project (or
% drag in the whole zip), set main.tex as the main document, and compile. No
% \graphicspath and no subdirectories -- the PDFs sit next to this file, so
% nothing has to be configured.
%
% This is a SELF-CONTAINED preview. To put the figure in a real paper, copy the
% figure* block below into it, add \usepackage{graphicx} and
% \usepackage{subcaption} to that paper's preamble, and point the three
% \includegraphics at wherever its figures live.
% ===========================================================================
\documentclass[10pt,twocolumn]{article}

% USENIX-shaped page: 7in text width, two 3.33in columns. Change or drop this
% line for a different venue -- the only thing it affects is how wide the float
% is allowed to be, and the panels are sized for a 7in full-width float.
\usepackage[letterpaper,textwidth=7in,columnsep=0.33in,textheight=9.25in]{geometry}

\usepackage{graphicx}
\usepackage{subcaption}
\usepackage{newtxtext}   % Times clone; drop it if the paper uses another face

\pagestyle{empty}
\begin{document}

% --------------------------------------------------------------- the figure
% Silicon is (a) and comes first: the hardware is the measurement and the
% simulator is the comparison.
%
% \includegraphics carries no width= on purpose. The panels are drawn at their
% final size with 9pt type, and 0.49\textwidth of a 7in textwidth is within 3%
% of their natural 3.35in, so leaving them unscaled keeps 9pt at 9pt. Scaling
% to fit a box would scale the type with it.
\begin{figure*}[t]
  \centering
  
\includegraphics{crossover-panel-legend}\\[-2pt]
  \begin{subfigure}[t]{0.49\textwidth}
    \centering
    \includegraphics{crossover-panel-a-m4}
    \caption{Apple M4, shipping silicon. $f^* = 29.3\%$.}
    \label{fig:xover-m4}
  \end{subfigure}\hfill
  \begin{subfigure}[t]{0.49\textwidth}
    \centering
    \includegraphics{crossover-panel-b-gem5}
    \caption{gem5 Neoverse-V2 FDP. $f^* = 60.0\%$.}
    \label{fig:xover-gem5}
  \end{subfigure}
  \caption{%
    \textbf{Where hand placement stops paying, on one microbenchmark and two
    machines.} A request is a public lookup lane of $L$ dependent loads and one
    AES-256-GCM call over 64 secret bytes; sweeping $L$ sweeps the secret
    fraction $f$ on the $x$ axis. Blanket \texttt{PSTATE.DIT} costs the public
    lane its load-value predictions, so it gets dearer as $f$ falls; Apple's
    bracket at the crypto call pays a fixed bill per call, so it gets cheaper.
    $f^*$ is where they cross, and it is the deployment decision. $y$ is IPC
    overhead against an unhardened build, each bracket arm measured against an
    instruction-matched NOP twin.
    %
    \emph{The two $y$ ranges differ, $+105\%$ against $+35\%$, and that is a
    result rather than a plotting choice}: the same bracket costs
    333--710\,cycles per request on the M4 against 206--351 in the model, over a
    request 3.6$\times$ shorter, so it is 106\% of a secret-heavy request on
    hardware and 18\% of one in simulation. The simulator is optimistic about
    hand placement by roughly a factor of two in secret fraction. The dashed
    curve is the simulator's own renamed switch design, which never crosses.%
  }
  \label{fig:xover}
\end{figure*}

\null   % the float needs a page to attach to; delete when pasting into a paper

\end{document}
